\section{Methodology}
\label{sec:methods}

\subsection{Study Cohorts}
Three non-synthetic clinical cohorts were used, all derived from SEER and TCGA-KIRC public registries. The \textbf{Multimodal Overlapping Cohort} contains TCGA-KIRC patients with aligned clinical, genomic, and CT radiomics data ($n = 126$, $M_1 = 18$, $M_0 = 108$, metastasis prevalence 14.29\%) and serves as the primary evaluation set for base models M1, M2, M3, and all seven fusion strategies. The \textbf{Genomic Extended Cohort} contains TCGA-KIRC patients with complete 19-gene expression signatures ($n = 418$, $M_1 = 54$, $M_0 = 364$, prevalence 12.92\%) and is used for secondary generalization evaluation of M2. The \textbf{SEER External Pre-training Cohort} ($n = 44,158$; $M_1 = 2,831$, $M_0 = 41,327$, prevalence 6.41\%) is used for external pre-training of M1 under zero-shot transfer. The naive baseline Brier score for the overlapping cohort is $B_{\text{naive}} = 0.1224$, corresponding to the constant prediction $p = 18/126 = 0.1429$.

\subsection{Base Model Architectures}
\textbf{M1 Clinical} is an XGBoost classifier trained on SEER ($n = 44,158$) using age, sex, tumor size, histological grade, and AJCC stage; it is evaluated zero-shot without fine-tuning on the TCGA-KIRC overlap ($n = 126$). \textbf{M2 Genomic} is a 5-fold out-of-fold XGBoost classifier trained with Synthetic Minority Over-sampling Technique (SMOTE) oversampling using a 19-gene clear cell RCC risk signature (\textit{VHL}, \textit{PBRM1}, \textit{SETD2}, \textit{BAP1}, \textit{KDM5C}); it is evaluated via 5-fold stratified cross-validation with SMOTE applied only within the training fold per iteration. \textbf{M3 CT Imaging} is a 5-fold out-of-fold Logistic Regression / linear SVM on 2,048-dimensional ResNet50 deep features extracted from pre-treatment abdominal CT scans, reduced via principal component analysis; the ResNet50 pathway captures high-attenuation tumor heterogeneity features corresponding to irregular tumor margins and necrotic core regions.

\subsection{Decision Fusion Strategies}
Seven decision-level fusion strategies were evaluated across three calibration regimes:
\begin{enumerate}
    \item \textbf{Fusion A (Simple Average):} Unweighted arithmetic mean $P = \frac{1}{3}(P_1 + P_2 + P_3)$.
    \item \textbf{Fusion B ($F_2$-Weighted Average, Primary Endpoint):} Weighted average $P = \frac{w_1 P_1 + w_2 P_2 + w_3 P_3}{w_1 + w_2 + w_3}$ with $w_1 = 0.4971, w_2 = 0.4918, w_3 = 0.5298$ (M1, M2, M3 validation $F_2$-scores).
    \item \textbf{Fusion C (Stacking Meta-Learner):} Out-of-fold logistic regression stacking trained on base-model probability predictions.
    \item \textbf{Fusion D (Cascade Max):} Rule-based maximum risk score $P = \max(P_1, P_2, P_3)$.
    \item \textbf{Fusion E (Bayesian Evidence Fusion, BEF):} Logit update $\text{logit}(P_{\text{fused}}) = \text{logit}(P_0) + \sum_m (\text{logit}(P_m) - \text{logit}(P_0))$, assuming conditional independence across modalities.
    \item \textbf{Fusion F (Dempster--Shafer Theory, DST):} Combines mass functions constructed from bounded belief and uncertainty intervals \cite{liu2017weighted, fontani2013dempster, huang2024deep}.
    \item \textbf{Fusion G (Optimal Transport, OT):} Combines prediction vectors by minimizing 1D Wasserstein distance to optimal target distributions \cite{courty2016optimal, kolouri2017optimal}.
\end{enumerate}

\subsection{Calibration Methodology}
Two calibration techniques were compared against the naive baseline. \textbf{Platt scaling} applies a parametric sigmoid transformation of model logits into calibrated probabilities fitted on a held-out validation fold; critically, Platt scaling is strictly rank-preserving per fold, with $\Delta \text{AUROC} = 0.00 \times 10^0$ across all 5 folds for all 3 base models, making AUROC comparisons across calibration regimes interpretable \cite{chen2016calibrating, maierhein2024metrics}. \textbf{Isotonic regression} is a non-parametric step function that introduced per-fold AUROC deviations ranging from 0.008 to 0.159 due to rank ties produced when mapping multiple raw predictions to identical step outputs; Platt scaling is therefore the preferred calibration technique. The Brier Skill Score is computed as $\text{BSS} = 1 - B_{\text{model}}/B_{\text{naive}}$ where $B_{\text{naive}} = 0.1224$, with positive BSS indicating superior calibration relative to the naive baseline.

\subsection{LLM-Powered Retrieval-Augmented Generation (RAG) Architecture}
\label{sec:rag_methods}

To extend the Calibration--Fusion Hazard diagnosis from decision-level classifier fusion to retrieval-augmented clinical decision support, we deployed a real LLM-powered RAG framework evaluated across three architectural tiers. All experiments used Qwen2.5-3B-Instruct (3.09 billion parameters, FP16 precision) as the generative backbone, loaded on an NVIDIA L40S GPU (48~GB VRAM), with PubMedBERT sentence-transformer embeddings (384-dimensional) for dense semantic retrieval.

\subsubsection{Clinical Knowledge Corpus}
A curated corpus of 10 ccRCC clinical guideline documents was assembled covering staging criteria, treatment algorithms, genomic biomarker interpretation, and distant metastasis risk assessment protocols. Each document was encoded into 384-dimensional dense embeddings via the sentence-transformer model, producing a pre-computed embedding matrix of shape $10 \times 384$.

\subsubsection{Multimodal Clinical Queries}
A standardized benchmark of 150 multimodal ccRCC clinical decision queries was generated, each encoding a realistic patient scenario with clinical staging text, genomic mutation profile, CT radiomic descriptor, and ground-truth distant metastasis label ($M_1 \in \{0, 1\}$), along with a ranked list of relevant guideline document identifiers for retrieval evaluation.

\subsubsection{RAG Architecture Variants}
Three RAG architectures of increasing complexity were evaluated:

\begin{enumerate}
    \item \textbf{Naive RAG:} A single dense semantic retriever (PubMedBERT cosine similarity) retrieves the top-$k$ ($k=3$) most similar guideline documents; the LLM generates a grounded clinical assessment from the concatenated retrieved context in a single forward pass.
    
    \item \textbf{Multimodal Hybrid RAG:} Three heterogeneous retrieval streams---sparse lexical (BM25), dense semantic (PubMedBERT cosine), and multimodal radiomic (feature-space similarity)---produce independent relevance scores that are fused via either uncalibrated normalization (Min-Max, Reciprocal Rank Fusion) or out-of-fold Platt/Isotonic calibrated log-odds fusion before the LLM generation step.
    
    \item \textbf{Agentic RAG:} A ReAct (Reasoning + Acting) tool-calling loop in which the LLM autonomously generates structured Thought $\to$ Action $\to$ Observation traces across up to 3 rounds, selecting from a tool inventory comprising a guideline retriever, a RenoFusion M1 risk calculator, and a genomic biomarker query engine. The agent decides when sufficient evidence has been gathered and issues a final clinical assessment via the \texttt{FINISH} action.
\end{enumerate}

\subsubsection{RAG Evaluation Metrics}
\textbf{Retrieval quality} was measured via Recall@$k$, Mean Reciprocal Rank (MRR), and Normalized Discounted Cumulative Gain (NDCG@3). \textbf{Generation quality} was measured via LLM-as-Judge Faithfulness (the Qwen model self-evaluates whether its generated answer is grounded in the retrieved context), Clinical Relevance (keyword-grounded coverage of ccRCC staging and risk terminology), Clinical Correctness (whether the generated assessment correctly identifies the ground-truth metastatic risk level), Hallucination Rate (proportion of ungrounded claims), and a composite RAGAS score (harmonic mean of Faithfulness and Relevance). \textbf{Calibration quality} was measured via Expected Calibration Error (ECE), Maximum Calibration Error (MCE), and Brier Score. All RAG experiments were evaluated under 5-fold stratified out-of-fold cross-validation with 2,000 paired bootstrap resamples.

\subsection{Statistical Analysis}
All base models and stacking meta-learners are evaluated via 5-fold stratified cross-validation, preserving the 14.29\% $M_1$ prevalence and ensuring no patient appears in both training and validation folds; SMOTE is applied only to the training fold per iteration. Paired empirical bootstrap (2,000 iterations) provides all AUROC CIs and fusion-versus-best-modality significance tests by resampling identical patient splits with replacement. The DeLong test is reported for methodological comparison with the paired bootstrap, recognizing that DeLong's asymptotic assumptions may be under-supported in small-positive cohorts ($N_{\text{pos}} = 18$). Decision curve analysis computes net benefit as $\text{NB}(p_t) = \frac{\text{TP}}{N} - \frac{\text{FP}}{N} \cdot \frac{p_t}{1 - p_t}$ across $p_t \in [0.01, 0.50]$. Expected and Maximum Calibration Error are computed across 8 probability bins. SHAP values are computed per base model to identify feature drivers: AJCC stage, tumor size, and histological grade cumulatively account for 84\% of clinical feature attribution; \textit{VHL}, \textit{PBRM1}, \textit{BAP1}, \textit{SETD2}, and \textit{KDM5C} dominate genomic attribution; and high-attenuation tumor heterogeneity ResNet50 features dominate imaging attribution. At $N = 150$ RAG queries, the paired bootstrap achieves 95.42\% statistical power ($\alpha = 0.05$, $d = 0.30$) with a minimum detectable effect size of $d = 0.2287$.
